% !TeX root = ./../main.tex
% Makra symboli stosowanych w pracy.
% Każdy nowy symbol dodaj tutaj oraz w wykazie chapters/symbols.tex.

% --- Informatyka kwantowa ---
\newcommand{\hilbert}{\mathcal{H}}                 % przestrzeń Hilberta
\newcommand{\identity}{\mathbf{1}}                 % operator identycznościowy
\newcommand{\prob}{\mathfrak{p}}                   % prawdopodobieństwo wyniku pomiaru
\newcommand{\channel}{\mathcal{E}}                 % kanał kwantowy

% --- Próbkowanie kwantowe i metoda PDU ---
\newcommand{\reffun}{\gamma}                       % funkcja referencyjna
\newcommand{\reffunexp}{\tilde{\gamma}}            % funkcja referencyjna wyznaczona eksperymentalnie
\newcommand{\pdu}{\varepsilon}                     % funkcja PDU (błąd)

% --- Kwantowe reprezentacje obrazów ---
\newcommand{\img}{I_m}                             % obraz (macierz intensywności)
\newcommand{\lpiqe}{\mathcal{L}}                   % operator kodowania LPIQE
\newcommand{\reconst}{\mathcal{R}}                 % operator rekonstrukcji obrazu

% --- Szum na poziomie generatorów ---
\newcommand{\liealg}[1]{\mathfrak{#1}}             % algebra Liego, np. \liealg{su}(2)
\newcommand{\genpert}{\Delta}                      % zaburzenie generatora

% --- Uczenie maszynowe ---
\newcommand{\dataset}{\mathcal{D}}
\newcommand{\trainset}{\dataset_{\mathrm{train}}}
\newcommand{\testset}{\dataset_{\mathrm{test}}}
\newcommand{\valset}{\dataset_{\mathrm{val}}}
\newcommand{\symloss}{\mathcal{L}_{\mathrm{loss}}} % funkcja straty (odróżniona od operatora LPIQE)
